\documentclass[11pt]{article}
\usepackage[margin=1in]{geometry}
\usepackage{xcolor}
\usepackage[breakable,listings]{tcolorbox}
\tcbuselibrary{listings}
\usepackage{hyperref}
\usepackage{enumitem}
\usepackage{listings}
\usepackage{verbatim}
\usepackage{amsmath}
\usepackage{graphicx}
\usepackage{booktabs}
\usepackage{float}

% Define colors
\definecolor{osaorange}{RGB}{220,115,38}
\definecolor{aiblue}{RGB}{30,144,255}
\definecolor{codebackground}{RGB}{248,248,248}

% Configure hyperref
\hypersetup{
    colorlinks=true,
    linkcolor=blue,
    urlcolor=blue,
    citecolor=blue
}

% Configure listings for code display
\lstset{
  basicstyle=\small\ttfamily,
  breaklines=true,
  breakatwhitespace=true,
  postbreak=\mbox{\textcolor{gray}{$\hookrightarrow$}\space},
  columns=flexible,
  keepspaces=true,
  showstringspaces=false
}

% Code environment using tcolorbox with listings
\newtcblisting{rcode}{
  colback=white,
  colframe=gray!50,
  boxrule=0.5pt,
  arc=0pt,
  left=3pt,
  right=3pt,
  top=3pt,
  bottom=3pt,
  width=\textwidth,
  breakable,
  listing only,
  listing options={
    basicstyle=\small\ttfamily,
    breaklines=true,
    breakatwhitespace=true,
    postbreak=\mbox{\textcolor{gray}{$\hookrightarrow$}\space},
    columns=flexible,
    keepspaces=true,
    showstringspaces=false
  }
}

\title{}
\author{}
\date{}

\begin{document}

% Header box
\begin{tcolorbox}[colback=osaorange,colframe=black,boxrule=2pt,arc=0pt,left=3pt,right=3pt,top=6pt,bottom=6pt,boxsep=2pt]
\centering
\textcolor{white}{\textbf{\Large H524 Introduction to Biostatistics}}\\
\textcolor{white}{\textbf{\Large Week 4 Lab: Confidence Intervals}}\\
\textcolor{white}{\textbf{\Large Fall 2025}}
\end{tcolorbox}

\vspace{8pt}

\noindent\textbf{Format:} Online\\
\textbf{Tools Required:} R, RStudio\\
\textbf{Estimated Time:} 90 minutes

\section{Learning Objectives}

By the end of this lab, you will be able to:
\begin{enumerate}
\item Calculate confidence intervals for a population mean using t-distribution
\item Use \texttt{t.test()} in R to compute CIs efficiently
\item Calculate and interpret confidence intervals for proportions
\item Understand how sample size, variability, and confidence level affect CI width
\item Determine required sample sizes for desired precision
\item Check normality assumptions using diagnostic plots
\item Visualize confidence intervals and their interpretation
\item \textcolor{aiblue}{\textbf{Use AI tools to help with CI calculations and interpretation}}
\end{enumerate}

\section{Part 1: Confidence Intervals for the Mean}

Confidence intervals provide a range of plausible values for an unknown population parameter. For a population mean with unknown variance, we use the t-distribution.

\subsection{Exercise 1.1: Manual CI Calculation}

\textbf{Scenario:} A study measures systolic blood pressure in 25 hypertensive patients.

\begin{rcode}
# Blood pressure data (mmHg)
bp_data <- c(128, 132, 125, 138, 130, 135, 127, 140, 129, 133,
             136, 126, 131, 134, 137, 142, 124, 139, 128, 135,
             130, 141, 127, 133, 136)

# Step 1: Calculate sample statistics
n <- length(bp_data)
xbar <- mean(bp_data)
s <- sd(bp_data)

cat("Sample size (n):", n, "\n")
cat("Sample mean:", round(xbar, 2), "mmHg\n")
cat("Sample SD:", round(s, 3), "mmHg\n\n")

# Step 2: Calculate standard error
se <- s / sqrt(n)
cat("Standard error:", round(se, 3), "mmHg\n\n")

# Step 3: Find critical value from t-distribution
df <- n - 1
conf_level <- 0.95
alpha <- 1 - conf_level
t_crit <- qt(1 - alpha/2, df)

cat("Degrees of freedom:", df, "\n")
cat("Confidence level:", conf_level*100, "%\n")
cat("t critical value:", round(t_crit, 4), "\n\n")

# Step 4: Calculate margin of error
margin_error <- t_crit * se
cat("Margin of error:", round(margin_error, 3), "mmHg\n\n")

# Step 5: Construct confidence interval
ci_lower <- xbar - margin_error
ci_upper <- xbar + margin_error

cat("95% Confidence Interval:\n")
cat("(", round(ci_lower, 2), ", ", round(ci_upper, 2), ") mmHg\n", sep="")
\end{rcode}

\textbf{Interpretation:} We are 95\% confident that the true mean systolic blood pressure for hypertensive patients is between the lower and upper bounds.

\vspace{0.5cm}

\textbf{Your turn:} What would happen to the CI if we used 90\% confidence instead? Calculate it below:

\vspace{2cm}

\subsection{Exercise 1.2: Using t.test() for CIs}

The \texttt{t.test()} function provides an easier way to calculate confidence intervals.

\begin{rcode}
# Same blood pressure data
# Use t.test() to get 95% CI
result <- t.test(bp_data, conf.level=0.95)

# Display the full output
print(result)

# Extract just the confidence interval
ci <- result$conf.int
cat("\n95% CI from t.test():", round(ci, 2), "\n")

# Compare with manual calculation
cat("Manual CI: (", round(ci_lower, 2), ", ", round(ci_upper, 2), ")\n", sep="")
cat("t.test CI: (", round(ci[1], 2), ", ", round(ci[2], 2), ")\n", sep="")
cat("These should match!\n")
\end{rcode}

\textbf{Advantages of t.test():}
\begin{itemize}
\item Simpler and faster
\item Less chance of calculation errors
\item Provides additional information (mean, t-statistic, p-value)
\item Easy to change confidence level
\end{itemize}

\subsection{Exercise 1.3: Different Confidence Levels}

Let's see how the confidence level affects the interval width.

\begin{rcode}
# Pain reduction scores from a clinical trial
pain_reduction <- c(42, 38, 51, 45, 39, 48, 44, 50, 43, 47,
                    41, 46, 49, 40, 52, 37, 45, 48, 42, 44)

# Calculate sample mean
mean_reduction <- mean(pain_reduction)
cat("Sample mean pain reduction:", round(mean_reduction, 2), "points\n\n")

# Calculate CIs for different confidence levels
conf_levels <- c(0.90, 0.95, 0.99)

for (conf in conf_levels) {
  result <- t.test(pain_reduction, conf.level=conf)
  ci <- result$conf.int
  width <- ci[2] - ci[1]

  cat(conf*100, "% CI: (", round(ci[1], 2), ", ", round(ci[2], 2),
      ")  Width = ", round(width, 2), "\n", sep="")
}

cat("\nNotice: Higher confidence = wider interval\n")
cat("Trade-off: More confidence means less precision\n")
\end{rcode}

\textbf{Question:} Why does higher confidence require a wider interval? Discuss with your partner.

\vspace{2cm}

\subsection{Exercise 1.4: Effect of Sample Size}

Sample size has a major impact on CI width. Let's demonstrate this.

\begin{rcode}
# Simulate the effect of sample size on CI width
set.seed(524)

# Population parameters (true values)
true_mean <- 100
true_sd <- 20

# Different sample sizes
sample_sizes <- c(10, 25, 50, 100, 200)

cat("Effect of Sample Size on 95% CI Width\n")
cat("(Population: mean=100, SD=20)\n\n")

results <- data.frame(n=sample_sizes,
                      CI_lower=NA, CI_upper=NA,
                      width=NA, SE=NA)

for (i in 1:length(sample_sizes)) {
  n <- sample_sizes[i]

  # Draw a sample
  sample_data <- rnorm(n, mean=true_mean, sd=true_sd)

  # Calculate CI
  result <- t.test(sample_data, conf.level=0.95)
  ci <- result$conf.int

  # Store results
  results$CI_lower[i] <- ci[1]
  results$CI_upper[i] <- ci[2]
  results$width[i] <- ci[2] - ci[1]
  results$SE[i] <- sd(sample_data) / sqrt(n)

  cat("n =", sprintf("%3d", n),
      "  SE =", sprintf("%5.2f", results$SE[i]),
      "  95% CI: (", round(ci[1], 1), ", ", round(ci[2], 1), ")",
      "  Width =", round(results$width[i], 1), "\n", sep="")
}

cat("\nObservation: Larger n gives narrower (more precise) CIs\n")
cat("SE decreases as 1/sqrt(n)\n")
\end{rcode}

\textbf{Question:} How much would you need to increase the sample size to cut the CI width in half?

\vspace{2cm}

\section{Part 2: Checking Assumptions}

The t-interval requires that the data come from an approximately normal distribution (or that $n$ is large enough for CLT).

\subsection{Exercise 2.1: Visual Checks for Normality}

\begin{rcode}
# Cholesterol data
cholesterol <- c(205, 198, 220, 185, 210, 195, 225, 200, 215, 192,
                 208, 196, 218, 203, 212, 189, 207, 201, 213, 199,
                 223, 191, 206, 214, 197)

# Create diagnostic plots
par(mfrow=c(1,3))  # 1 row, 3 columns

# 1. Histogram
hist(cholesterol, breaks=10, col="lightblue", border="white",
     main="Histogram", xlab="Cholesterol (mg/dL)",
     freq=FALSE)  # density scale

# Overlay normal curve
xfit <- seq(min(cholesterol), max(cholesterol), length=100)
yfit <- dnorm(xfit, mean=mean(cholesterol), sd=sd(cholesterol))
lines(xfit, yfit, col="red", lwd=2)

# 2. Boxplot (check for outliers)
boxplot(cholesterol, col="lightgreen",
        main="Boxplot", ylab="Cholesterol (mg/dL)")

# 3. Q-Q plot (most important!)
qqnorm(cholesterol, main="Q-Q Plot", pch=19, col="blue")
qqline(cholesterol, col="red", lwd=2)

par(mfrow=c(1,1))  # Reset layout

# Interpretation guide
cat("\nAssessing Normality:\n")
cat("- Histogram: Should be roughly bell-shaped\n")
cat("- Boxplot: Check for extreme outliers\n")
cat("- Q-Q plot: Points should fall on the line\n")
cat("  (slight deviations at ends are OK)\n")
\end{rcode}

\textbf{Assessment:} Based on these plots, do the cholesterol data appear approximately normal?

\vspace{2cm}

\subsection{Exercise 2.2: Shapiro-Wilk Test}

A formal statistical test for normality.

\begin{rcode}
# Shapiro-Wilk test
shapiro_result <- shapiro.test(cholesterol)

cat("Shapiro-Wilk Test for Normality\n")
cat("H0: Data come from a normal distribution\n\n")
cat("Test statistic W:", round(shapiro_result$statistic, 4), "\n")
cat("P-value:", round(shapiro_result$p.value, 4), "\n\n")

if (shapiro_result$p.value > 0.05) {
  cat("Conclusion: No evidence against normality (p > 0.05)\n")
  cat("Safe to use t-interval\n")
} else {
  cat("Conclusion: Evidence of non-normality (p < 0.05)\n")
  cat("Consider: transformation or larger sample size\n")
}
\end{rcode}

\textbf{Note:} With large samples, the test becomes very sensitive and may reject normality for trivial departures. Visual inspection is often more useful.

\subsection{Exercise 2.3: Robustness of t-Interval}

The t-interval is fairly robust to mild departures from normality, especially with larger samples.

\begin{rcode}
# Demonstrate robustness with slightly skewed data
set.seed(123)

# Generate right-skewed data (exponential)
skewed_data <- rexp(30, rate=0.1)

# Check normality
par(mfrow=c(1,2))
hist(skewed_data, breaks=10, col="lightcoral", border="white",
     main="Skewed Data", xlab="Value")
qqnorm(skewed_data, main="Q-Q Plot", pch=19)
qqline(skewed_data, col="red", lwd=2)
par(mfrow=c(1,1))

# Calculate CI anyway (n=30 is borderline)
result <- t.test(skewed_data, conf.level=0.95)
cat("\n95% CI for skewed data:", round(result$conf.int, 2), "\n")
cat("Sample mean:", round(mean(skewed_data), 2), "\n")

# Shapiro test
cat("\nShapiro test p-value:", round(shapiro.test(skewed_data)$p.value, 4), "\n")
cat("\nWith n=30, t-interval is reasonably robust\n")
cat("For severe skewness or n<15, consider alternatives\n")
\end{rcode}

\section{Part 3: Confidence Intervals for Proportions}

When the outcome is binary (success/failure), we estimate a population proportion.

\subsection{Exercise 3.1: Basic CI for Proportion}

\begin{rcode}
# Vaccine efficacy study
# 174 out of 200 developed immunity
n <- 200
x <- 174  # Number of successes
p_hat <- x / n

cat("Sample proportion:", p_hat, "\n")
cat("Percentage:", p_hat*100, "%\n\n")

# Check conditions for normal approximation
cat("Checking conditions:\n")
cat("n*p_hat =", n*p_hat, "(should be >= 10)\n")
cat("n*(1-p_hat) =", n*(1-p_hat), "(should be >= 10)\n")

if (n*p_hat >= 10 && n*(1-p_hat) >= 10) {
  cat("Conditions met! Can use normal approximation.\n\n")
} else {
  cat("Conditions NOT met. Use exact methods.\n\n")
}

# Calculate standard error for proportion
se_prop <- sqrt(p_hat * (1 - p_hat) / n)
cat("Standard error:", round(se_prop, 4), "\n\n")

# 95% CI using normal approximation
z_crit <- qnorm(0.975)  # 1.96
margin <- z_crit * se_prop

ci_lower <- p_hat - margin
ci_upper <- p_hat + margin

cat("95% Confidence Interval for proportion:\n")
cat("(", round(ci_lower, 4), ", ", round(ci_upper, 4), ")\n", sep="")
cat("Or: (", round(ci_lower*100, 2), "%, ", round(ci_upper*100, 2), "%)\n", sep="")

# Alternative: use prop.test()
prop_result <- prop.test(x, n, conf.level=0.95, correct=FALSE)
cat("\nUsing prop.test():\n")
print(prop_result$conf.int)
\end{rcode}

\textbf{Interpretation:} We are 95\% confident that the true vaccine efficacy is between the lower and upper bounds.

\subsection{Exercise 3.2: Disease Prevalence}

\begin{rcode}
# Diabetes screening study
# 65 out of 500 adults test positive
n_screen <- 500
positive <- 65
prevalence <- positive / n_screen

cat("Sample prevalence:", prevalence, "\n")
cat("As percentage:", prevalence*100, "%\n\n")

# Calculate 95% CI
se_prev <- sqrt(prevalence * (1 - prevalence) / n_screen)
z <- 1.96

ci_prev_lower <- prevalence - z * se_prev
ci_prev_upper <- prevalence + z * se_prev

cat("95% CI for prevalence:\n")
cat("(", round(ci_prev_lower*100, 2), "%, ",
    round(ci_prev_upper*100, 2), "%)\n", sep="")

# Using prop.test
prop.test(positive, n_screen, conf.level=0.95, correct=FALSE)
\end{rcode}

\textbf{Your turn:} If 150 out of 400 patients responded to treatment, calculate the 95\% CI for the response rate.

\vspace{3cm}

\section{Part 4: Sample Size Determination}

Before conducting a study, we need to determine how many subjects to recruit.

\subsection{Exercise 4.1: Sample Size for Estimating a Mean}

\begin{rcode}
# Function to calculate required sample size
sample_size_mean <- function(sigma, margin_error, conf_level=0.95) {
  alpha <- 1 - conf_level
  z_crit <- qnorm(1 - alpha/2)

  n <- (z_crit * sigma / margin_error)^2

  return(ceiling(n))  # Always round up!
}

# Example: Cholesterol study
# Want margin of error = 5 mg/dL
# Literature suggests SD = 40 mg/dL

sigma_est <- 40
me_desired <- 5

n_required <- sample_size_mean(sigma_est, me_desired, 0.95)

cat("Sample Size Calculation\n")
cat("=======================\n")
cat("Estimated population SD:", sigma_est, "mg/dL\n")
cat("Desired margin of error:", me_desired, "mg/dL\n")
cat("Confidence level: 95%\n\n")
cat("Required sample size:", n_required, "\n")
\end{rcode}

\subsection{Exercise 4.2: How Margin of Error Affects Sample Size}

\begin{rcode}
# Blood pressure study planning
# Assume SD = 15 mmHg

sigma_bp <- 15

# Try different margins of error
margins <- c(2, 3, 4, 5, 10)

cat("Blood Pressure Study Planning\n")
cat("(Estimated SD = 15 mmHg, 95% confidence)\n\n")
cat("Desired ME    Required n\n")
cat("----------    ----------\n")

for (me in margins) {
  n <- sample_size_mean(sigma_bp, me, 0.95)
  cat(sprintf("%5d mmHg", me), "       ", sprintf("%4d", n), "\n")
}

cat("\nKey insight: To halve ME, need 4x the sample size!\n")
\end{rcode}

\subsection{Exercise 4.3: Sample Size for Proportion}

\begin{rcode}
# Function for proportion sample size
sample_size_prop <- function(p, margin_error, conf_level=0.95) {
  alpha <- 1 - conf_level
  z_crit <- qnorm(1 - alpha/2)

  n <- (z_crit / margin_error)^2 * p * (1 - p)

  return(ceiling(n))
}

# Prevalence study
# Want ME = 0.03 (3 percentage points)

me_prev <- 0.03

# Scenario 1: No prior information (use p=0.5 - conservative)
n_conservative <- sample_size_prop(0.5, me_prev, 0.95)

# Scenario 2: Prior studies suggest p = 0.15
n_informed <- sample_size_prop(0.15, me_prev, 0.95)

cat("Sample Size for Prevalence Study\n")
cat("(Desired ME = 3%, 95% confidence)\n\n")
cat("Conservative (p=0.5):", n_conservative, "subjects\n")
cat("Informed (p=0.15):", n_informed, "subjects\n\n")
cat("Savings from prior information:", n_conservative - n_informed, "subjects\n")

# Show how p affects required n
cat("\nEffect of assumed proportion:\n")
proportions <- c(0.1, 0.2, 0.3, 0.4, 0.5)
for (p in proportions) {
  n <- sample_size_prop(p, me_prev, 0.95)
  cat("p =", p, " -> n =", n, "\n")
}
cat("\np=0.5 gives maximum n (most conservative)\n")
\end{rcode}

\section{Part 5: Visualizing Confidence Intervals}

Understanding what confidence intervals mean through simulation.

\subsection{Exercise 5.1: The Meaning of 95\% Confidence}

\begin{rcode}
# Simulate many samples and their CIs
set.seed(524)

true_mean <- 100   # True population mean (usually unknown!)
true_sd <- 15      # True population SD
sample_size <- 25
num_samples <- 20  # Draw 20 different samples

# Storage
ci_results <- data.frame(sample_num=1:num_samples,
                         lower=NA, upper=NA,
                         sample_mean=NA,
                         contains_true_mean=NA)

# Draw samples and calculate CIs
for (i in 1:num_samples) {
  # Draw a random sample
  sample_data <- rnorm(sample_size, mean=true_mean, sd=true_sd)

  # Calculate CI
  result <- t.test(sample_data, conf.level=0.95)
  ci <- result$conf.int

  # Store results
  ci_results$lower[i] <- ci[1]
  ci_results$upper[i] <- ci[2]
  ci_results$sample_mean[i] <- mean(sample_data)
  ci_results$contains_true_mean[i] <- (ci[1] <= true_mean) & (true_mean <= ci[2])
}

# How many CIs contain the true mean?
num_contain <- sum(ci_results$contains_true_mean)
cat("Out of", num_samples, "confidence intervals,",
    num_contain, "contain the true mean\n")
cat("Percentage:", round(num_contain/num_samples*100, 1), "%\n")
cat("Expected: ~95%\n")
\end{rcode}

\subsection{Exercise 5.2: Plotting Multiple CIs}

\begin{rcode}
# Visualize the confidence intervals
plot(NULL, xlim=c(85, 115), ylim=c(0, num_samples+1),
     xlab="Value", ylab="Sample Number",
     main="20 Confidence Intervals (95% level)")

# Add vertical line at true mean
abline(v=true_mean, col="blue", lwd=2, lty=1)

# Draw each CI
for (i in 1:num_samples) {
  # Color: green if contains true mean, red if doesn't
  color <- ifelse(ci_results$contains_true_mean[i],
                  "darkgreen", "red")

  # Draw horizontal line for CI
  segments(ci_results$lower[i], i,
           ci_results$upper[i], i,
           col=color, lwd=2)

  # Add point at sample mean
  points(ci_results$sample_mean[i], i,
         pch=19, col=color, cex=0.8)
}

# Add legend
legend("topleft",
       legend=c("True mean", "Contains true mean", "Misses true mean"),
       col=c("blue", "darkgreen", "red"),
       lty=c(1, 1, 1), lwd=c(2, 2, 2))

text(115, num_samples,
     paste(num_contain, "out of", num_samples),
     pos=2, col="darkgreen", font=2)
\end{rcode}

\textbf{Key insight:} The confidence level (95\%) refers to the long-run frequency. In repeated sampling, 95\% of intervals will contain the true parameter.

\section{Practice Problems}

\subsection{Problem 1: Clinical Trial}

A new drug for lowering cholesterol is tested on 30 patients. The mean reduction is 28 mg/dL with SD = 12 mg/dL.

\begin{enumerate}[label=\alph*)]
\item Calculate a 95\% confidence interval for the true mean reduction
\item Calculate a 99\% confidence interval
\item Interpret both intervals in context
\item If the goal was a reduction of at least 25 mg/dL, what does your CI tell you?
\end{enumerate}

\vspace{3cm}

\subsection{Problem 2: Adverse Events}

In a safety trial, 8 out of 150 patients experienced adverse events.

\begin{enumerate}[label=\alph*)]
\item Check whether normal approximation is appropriate
\item Calculate a 95\% CI for the true proportion of adverse events
\item Express your answer as a percentage
\end{enumerate}

\vspace{3cm}

\subsection{Problem 3: Study Planning}

You're planning a study to estimate mean body temperature with 95\% confidence.

\begin{enumerate}[label=\alph*)]
\item If you want margin of error = 0.2°F and you estimate SD = 0.7°F, how many subjects do you need?
\item What if you want ME = 0.1°F instead?
\item How does doubling precision affect required sample size?
\end{enumerate}

\vspace{3cm}

\subsection{Problem 4: Assumption Checking}

Load this dataset and check normality assumptions:

\begin{rcode}
reaction_times <- c(245, 267, 289, 234, 256, 278, 298, 241,
                    263, 287, 312, 338, 229, 251, 273)

# Create histogram, boxplot, and Q-Q plot
# Run Shapiro-Wilk test
# Assess whether t-interval is appropriate
\end{rcode}

\vspace{3cm}

\section{\textcolor{aiblue}{AI-Enhanced Learning Tips}}

\begin{tcolorbox}[colback=aiblue!10,colframe=aiblue,title=Using AI Tools for Confidence Intervals]

\textbf{Good prompts for AI assistance:}
\begin{itemize}
\item ``Calculate a 95\% CI for mean given n=20, mean=45, SD=8''
\item ``Show me step-by-step how to compute margin of error''
\item ``What sample size do I need for ME=3 if SD=15 at 95\% confidence?''
\item ``Explain why we use t-distribution instead of normal distribution''
\item ``Create R code to check normality with histogram and Q-Q plot''
\item ``How do I interpret a confidence interval for a clinical audience?''
\end{itemize}

\textbf{AI can help with:}
\begin{itemize}
\item Step-by-step CI calculations
\item Explaining when to use different methods (t vs Z, mean vs proportion)
\item Generating diagnostic plots
\item Sample size calculations
\item Writing proper interpretations
\item Debugging R code errors
\end{itemize}

\textbf{Verification checklist:}
\begin{itemize}
\item Check that AI used correct formula (t for mean with unknown SD)
\item Verify degrees of freedom are correct (n-1)
\item Ensure critical value matches confidence level
\item Confirm standard error calculation
\item Cross-check with \texttt{t.test()} output
\end{itemize}

\textbf{Practice prompt:}
``I have cholesterol data from 25 patients: mean=205, SD=42. Calculate a 95\% CI step-by-step, explain each step, and show how to verify it in R using t.test().''
\end{tcolorbox}

\section{Key Takeaways}

\begin{enumerate}
\item \textbf{t-Interval for mean:} $\bar{x} \pm t_{\alpha/2, n-1} \times \frac{s}{\sqrt{n}}$
  \begin{itemize}
  \item Use when population SD is unknown (almost always!)
  \item Requires approximate normality or large n ($\geq$ 30)
  \end{itemize}

\item \textbf{CI for proportion:} $\hat{p} \pm z_{\alpha/2} \times \sqrt{\frac{\hat{p}(1-\hat{p})}{n}}$
  \begin{itemize}
  \item Requires $n\hat{p} \geq 10$ and $n(1-\hat{p}) \geq 10$
  \item Uses normal (Z) distribution
  \end{itemize}

\item \textbf{Factors affecting CI width:}
  \begin{itemize}
  \item Larger n $\to$ narrower CI (more precise)
  \item Larger s $\to$ wider CI (less precise)
  \item Higher confidence $\to$ wider CI (more certain)
  \end{itemize}

\item \textbf{Interpretation:} ``We are 95\% confident that the true population parameter is between [lower, upper]''

\item \textbf{Sample size:} $n = \left(\frac{z \times \sigma}{\text{ME}}\right)^2$ (always round up!)

\item \textbf{R functions:}
  \begin{itemize}
  \item \texttt{t.test(data, conf.level=0.95)} - CI for mean
  \item \texttt{prop.test(x, n, conf.level=0.95)} - CI for proportion
  \item \texttt{qt(p, df)} - t critical values
  \item \texttt{qnorm(p)} - Z critical values
  \end{itemize}
\end{enumerate}

\section{Saving Your Work}

\begin{rcode}
# Save your workspace
save.image("week4_lab.RData")

# To load later:
# load("week4_lab.RData")

# Save your R script with header:
# Name: Your Name
# Date: Today's Date
# Lab: Week 4 - Confidence Intervals
\end{rcode}

\end{document}
